\chapter{Two-Sided Causality as a Time-Symmetric Projection from the Timeless Information Space}

\begin{qrabstract}
This chapter formulates a QR-compatible description of \emph{two-sided causality}
as a time-symmetric projection from the timeless information space~$\mathcal{I}$.
The orientation of the arrow of time is not fundamental but emerges as a property
of the projection operator. To this end, we define the forward-oriented
operator family $\mathcal{O}^{(n)}_{\rightarrow}$ and its time-symmetric
extension $\mathcal{O}^{(n)}_{\pm}$. We show how even/odd derivative orders
modulate the causal structure, provide a scale-dependent selection criterion for $n$
via the smoothness parameter $\Xi_n$, and summarize the regimes (\emph{Expansion, Rotation,
Quantum}) in a table. Finally, we outline potential experimental signatures.
\end{qrabstract}

\section{Axiomatic Embedding}
The core QR concept is: the observed spacetime $M^{(4)}$ is the projection
\(\pi:\mathcal{I}\to M^{(4)}\) of a timeless information space~\(\mathcal{I}\).
The observed \emph{time} and hence causal structures are emergent properties
of the projection process. The universal QR observable is modeled as
\begin{equation}
\label{eq:qr-core}
\mathcal{O}(x^\mu,t)\;=\;\biggl(\frac{\partial^{n}}{\partial t^{n}}\log I_\nu(t)\biggr)\cdot
\mathcal{W}(\rho_{\mathrm{bar}},\|\psi\|,r)\cdot \mathcal{F}(r)
\end{equation}
with projection order \(n\in\mathbb{N}\) (physically: degree of temporal
information dynamics), structural weighting \(\mathcal{W}\), and fractal amplification
\(\mathcal{F}(r)=(r/\ell_\pi)^{\epsilon}\), where \(d_f=4+\epsilon\) denotes the effective fractal dimension
(\(0<\epsilon\ll1\)).

\section{Forward-Oriented and Symmetric Projection}
\subsection{Oriented Projection}
The \emph{forward-oriented} projection (classical arrow of time) is given by
\begin{equation}
\label{eq:forward}
\mathcal{O}^{(n)}_{\rightarrow}(t)\;=\;\frac{\partial^{n}}{\partial t^{n}}\log I_\nu(t)
\end{equation}
and produces, on large scales (\(r\gg \ell_\pi\)), the familiar empirically stable
causal order (past \(\to\) future).

\subsection{Time-Symmetric Projection}
Two-sided causality is described by a \emph{time-symmetric} operator choice:
\begin{equation}
\label{eq:symmetric}
\mathcal{O}^{(n)}_{\pm}(t)\;=\;\frac{1}{2}\Biggl[
\frac{\partial^{n}}{\partial t^{n}} + (-1)^{n}\frac{\partial^{n}}{\partial (-t)^{n}}
\Biggr]\log I_\nu(t)
\;=\;\begin{cases}
\displaystyle \frac{\partial^{n}}{\partial t^{n}}\log I_\nu(t), & n\ \text{even},\\[1.0ex]
\displaystyle 0, & n\ \text{odd},
\end{cases}
\end{equation}
where the last equality follows from the parity property.\footnote{%
This representation serves as a canonical reference form. Real systems with finite
coherence, finite boundaries, and \(\mathcal{W}\neq \mathrm{const.}\) generally yield
\(\mathcal{O}^{(n)}_{\pm}\neq 0\) for odd \(n\), but suppressed.}
Physically:
\begin{itemize}
  \item For \emph{even} \(n\), forward and backward contributions add constructively
  \(\Rightarrow\) enhanced sensitivity to two-sided (time-entangled) correlations.
  \item For \emph{odd} \(n\), contributions cancel partially
  \(\Rightarrow\) effective indefiniteness in temporal ordering (indefinite order).
\end{itemize}

\section{Scale Rule: Choice of Derivative Order \(n\)}
\subsection{Fractal Amplification and Smoothness}
The fractal amplification \(\mathcal{F}(r)=(r/\ell_\pi)^{\epsilon}\) acts as a scale-dependent
\emph{geometry amplifier}: the larger \(r\), the stronger the geometry contributes;
the smaller \(r\) (or the shorter the characteristic time \(\tau\)), the more higher
derivatives must explicitly capture the rapid information dynamics.

We define the dimensionless smoothness parameter
\begin{equation}
\label{eq:xi}
\Xi_n(\tau)\;:=\;\tau^{\,n}\,
\frac{\bigl|\partial_t^{\,n+1}\log I_\nu\bigr|}{\bigl|\partial_t^{\,n}\log I_\nu\bigr|}\,.
\end{equation}
\textbf{Selection rule:} Choose the smallest \(n\) such that
\begin{equation}
\label{eq:rule}
\Xi_n(\tau)\;\lesssim\;\epsilon
\end{equation}
holds. This ensures that the truncation error of the next-higher term is capped
by the fractal “roughness” \(\epsilon\) (parameter-free).

\subsection{Regimes and Standard Assignments}
The following table summarizes practical assignments. It is compatible
with the fractal scale ladder \(\lambda_k=\ell_\pi\,\varphi^{k}\) and established QR applications.

\begin{table}[h]
\centering
\caption{Scale-dependent assignment of derivative order and causal structure.}
\label{tab:regimes}
\begin{tabular}{@{}lcccc@{}}
\toprule
Regime & Typical scale \(L\) & Projection & Dominant \(n\) & Causal structure \\
\midrule
Cosmic expansion (background) & \(\gtrsim \mathrm{Gpc}\) & \(\mathcal{O}^{(1)}_{\rightarrow}\) & 1 & forward-dominated \\
Acoustics/structure (BAO/CMB) & \(10^2\,\mathrm{Mpc}\) & \(\mathcal{O}^{(2)}_{\pm}\) & 2 & partly symmetric \\
Galactic rotation & \(10\text{–}30\,\mathrm{kpc}\) & \(\mathcal{O}^{(2)}_{\pm}\) & 2 & partly symmetric \\
Solar system/GR limit & \(\ll 1\,\mathrm{pc}\) & \(\mathcal{O}^{(1)}_{\rightarrow}\) & 1 & classical/stable \\
Quantum/lab scale & \(\ll \mathrm{sub\text{-}pc}\) & \(\mathcal{O}^{(4)}_{\pm}\) & 4 & (nearly) symmetric \\
\bottomrule
\end{tabular}
\end{table}

\section{Experimental Signatures}
The time-symmetric projection yields clear, testable consequences:
\begin{enumerate}
  \item \textbf{Indefinite causal order:} For \(n\ge 3\), process sequences are not
  uniquely defined (quantum switch-like phenomena).
  \item \textbf{Delayed choice:} Projection is set by coherence conditions in information space;
  statistical effects correspond to delayed measurement choice.
  \item \textbf{Scale-dependent arrow-of-time stability:} Transition from forward-dominated
  causality (\(n=1\)) to (nearly) two-sided correlations (\(n=4\)) with decreasing
  \(\tau\) or \(L\).
\end{enumerate}

\section{Remark on Consistency}
Equations \eqref{eq:qr-core}–\eqref{eq:rule} are consistent with the QR axioms:
(i) Time emerges only through projection (no contradiction with timeless \(\mathcal{I}\));
(ii) \(\mathcal{F}(r)\) respects the fractal effective dimension \(d_f=4+\epsilon\);
(iii) the choice of \(n\) is \emph{not} free, but operationally fixed by the smoothness criterion
\(\Xi_n(\tau)\lesssim \epsilon\).

\section{Quick Guide (Practical)}
\begin{itemize}
  \item \textbf{Given:} Scale \(L\) (or time \(\tau\)), estimate for \(\epsilon\) (typically \(\ll 1\)).
  \item \textbf{Step 1:} Determine \(\tau\) (orbital, growth, coherence time).
  \item \textbf{Step 2:} Check \(\Xi_1(\tau)\); if \(\Xi_1\le\epsilon\Rightarrow n=1\),
  otherwise increase \(n\) until \(\Xi_n\le\epsilon\).
  \item \textbf{Step 3:} Choose projection type: large scale \(\Rightarrow \mathcal{O}^{(n)}_{\rightarrow}\),
  small scale/quantum \(\Rightarrow \mathcal{O}^{(n)}_{\pm}\).
\end{itemize}